\section{Discussion}
\label{sec:discussion}

\subsection{Quantization as an Implicit Regularizer}
The central theoretical insight established by this study is that \emph{uniform integer post-training quantization functions as an implicit regularization operator}. In high-capacity floating-point models, gradient optimization fits extremely subtle, high-frequency decision boundaries that maximize empirical separation on source training telemetry. When evaluated on target captures collected under different network topologies, these subtle boundaries transform into catastrophic failure modes.

By restricting weight tensors $\mathbf{W}$ to discrete steps $\Delta W = S_W$, PTQ acts analogously to weight clipping and stochastic rounding. Non-causal features whose marginal contribution to the log-odds is smaller than the quantization bin width are effectively neutralized. This explains the marked suppression of false alarms on UNSW-NB15 (reducing false positives from $5,612$ to $3,090$, a $44.9\%$ reduction) while simultaneously improving malicious detection count by $+310$ flows.

\subsection{The Reality of Standalone Zero-Shot Edge NIDS}
While INT8 significantly mitigates collapse ($+1.09\%$ accuracy gain, $\chi^2 = 1,871.69, p < 10^{-15}$), our findings underscore a sobering reality for practical edge cybersecurity: \textbf{Zero-shot cross-domain ML-NIDS remains unviable for standalone production deployment without domain adaptation}.

Even with quantization mitigation, the absolute attack recall across UNSW-NB15 is only $2.54\%$. For four out of nine attack families (Analysis, Backdoors, Reconnaissance, and Shellcode), zero-shot recall is exactly $0.00\%$. A system that misses $97.46\%$ of operational attacks cannot serve as a primary security perimeter. Consequently, we argue that the academic practice of evaluating ML-NIDS exclusively on holdout splits of single datasets (such as reporting $>99\%$ accuracy on CIC-IDS-2017) fosters a dangerously inflated sense of defensive capability.

\subsection{Methodological Strengths & Ethical Safeguards}
In response to widespread reproducibility and evaluation pitfalls documented by Arp et al.~\cite{arp2022dos}, ZeroEdge-NIDS embodies several rigorous experimental safeguards:
\begin{enumerate}
    \item \textbf{No Target Snooping:} Feature transformations, logarithmic scales, and min-max normalization parameters were strictly frozen on the CIC-IDS-2017 training set and never adapted to UNSW-NB15.
    \item \textbf{Zero Synthetic Data Fabrication:} All reported evaluations utilize real, unaugmented network traffic captures from verified public corpora.
    \item \textbf{Explicit Architectural Transparency:} We strictly segregate empirical timings measured on host hardware from analytical projections computed for the virtual ESP32 target.
    \item \textbf{Evaluation Dependency Gating:} The Interactive Experiment Lab prevents misleading metric reporting by requiring validated ground truth labels before calculating diagnostic classification scores.
\end{enumerate}

\subsection{Threats to Validity & Limitations}
\begin{itemize}
    \item \textbf{Simulation Boundaries:} The ESP32 latency, memory, and power figures reported in Section~\ref{sec:results} are mathematical projections derived from counted MAC operations and documented datasheet constants. Physical silicon measurements may introduce cache misses, memory bus contention, and peripheral interrupt jitter.
    \item \textbf{Coarse Binary Classification:} This study models network intrusion detection as a binary classification task (benign vs. malicious). While binary triage is typical for edge microcontrollers acting as lightweight packet filters, fine-grained multi-class attribution across specific CVEs remains unaddressed.
    \item \textbf{Batch Size Discrepancies:} Host CPU benchmarks utilize batch sizes of one flow to reflect true stream-based edge arrival, but CPU instruction caches and branch predictors still favor host execution over bare-metal microcontroller pipelines.
\end{itemize}
